\documentclass[11pt]{article}
\usepackage[margin=1in]{geometry}
\usepackage{graphicx}
\usepackage{booktabs}
\usepackage{amsmath, amssymb}
\usepackage{siunitx}
\usepackage{hyperref}

\title{Optech Raman AI: Résultats de pipeline}
\author{TTM (Polytechnique Montréal)}
\date{\today}

\begin{document}
\maketitle

\section{Résumé}
Ce document présente les étapes du pipeline (conversion, prétraitement, extraction de features, modélisation) et un aperçu des métriques.

\section{Pipeline}
\begin{itemize}
\item Conversion CSV $\rightarrow$ Parquet
\item Prétraitement (Savitzky--Golay, baseline ALS, normalisation à l'aire)
\item Features physiques (pics, positions, aires)
\item Modélisation bayésienne multi-sorties (exemple)
\end{itemize}

\section{Métriques}
\input{metrics_table.tex}

\bibliographystyle{ieeetr}
\bibliography{references}
\end{document}
